\subsection{Messwerte der Verifikationsversuche}\label{sec:AnhangMesswerte}
Die folgenden Tabellen enthalten die Rohdaten der in \cref{tab:pruefplan} aufgeführten
Verifikationsversuche. Die Spalte \emph{RUN} bezeichnet den fortlaufenden Durchlauf
einer Messreihe.

Die Messprotokolle wurden in Skalenteilen der Messuhr geführt. Bei einem
Skalenteilungswert von \qty{0.01}{\milli\meter} entspricht ein Skalenteil
\qty{10}{\micro\meter}; alle Werte sind hier bereits in Millimeter umgerechnet.
Der vor jeder Messreihe eingestellte Normierungswert dient ausschließlich dazu, den
Tastkopf in die Mitte des Messbereichs vorzuspannen, damit Auslenkungen in beide
Richtungen erfasst werden können. Er wurde von Hand eingestellt und stellt keine
rückführbare Referenz dar. Der absolute Ablesewert einer Messreihe ist damit für sich
genommen nicht aussagekräftig und zwischen Messreihen nicht vergleichbar.

Bewertet wird deshalb die Streuung innerhalb einer Messreihe. Dazu wird jeder
Ablesewert $x_i$ auf den Mittelwert $\bar{x}$ der eigenen Messreihe bezogen:
%
\begin{equation}
  \Delta_i = x_i - \bar{x},
  \qquad
  \bar{x} = \frac{1}{n}\sum_{i=1}^{n} x_i .
  \label{eq:abweichung}
\end{equation}
%
Da $\bar{x}$ den Normierungswert der jeweiligen Messreihe enthält, fällt dieser bei
der Differenzbildung heraus. Die so gebildeten Abweichungen $\Delta_i$ besitzen für
jede Messreihe den gemeinsamen Bezugspunkt null und sind damit unmittelbar
vergleichbar, unabhängig davon, auf welchen Wert die Messuhr zuvor vorgespannt wurde.
Als Kennwerte werden die empirische Standardabweichung $s$ und die Spanne
$R = x_{\max} - x_{\min}$ angegeben. Beide sind ebenfalls unabhängig vom
Normierungswert und beschreiben die Wiederholpräzision des jeweiligen Subsystems.

%----------------------------------------------------------------------
% T-01
%----------------------------------------------------------------------
\begin{table}[htbp]
  \centering
  \caption[Messwerte T-01: Sperrklinken (Gearlinear-System)]{Messwerte des Versuchs
    T-01 zur Positionstreue der Sperrklinkenschlitten. Setup: Messuhr an der
    Linearstange bzw.\ am Endeffektor, je zehn Durchführungen pro Messstelle.
    Die Messung an der Linearstange erfolgte direkt auf Stahl. $x$ bezeichnet den
    Ablesewert, $\Delta$ die Abweichung vom Mittelwert der jeweiligen Messreihe
    nach \cref{eq:abweichung}.}
  \label{tab:mess_T01}
  \small
  \begin{tabular}{@{}c S[table-format=1.2] S[table-format=+1.3] S[table-format=1.2] S[table-format=+1.3]@{}}
    \toprule
    & \multicolumn{2}{c}{\textbf{Linearstange}}
    & \multicolumn{2}{c}{\textbf{Endeffektor}} \\
    \cmidrule(lr){2-3}\cmidrule(l){4-5}
    \textbf{Durchlauf} & {$x$} & {$\Delta$} & {$x$} & {$\Delta$} \\
    & {\unit{\milli\meter}} & {\unit{\milli\meter}}
    & {\unit{\milli\meter}} & {\unit{\milli\meter}} \\
    \midrule
     1 & 0.35  & +0.007 & 0.58  & +0.010 \\
     2 & 0.35  & +0.007 & 0.59  & +0.020 \\
     3 & 0.34  & -0.003 & 0.55  & -0.021 \\
     4 & 0.34  & -0.003 & 0.55  & -0.021 \\
     5 & 0.32  & -0.023 & 0.57  & -0.001 \\
     6 & 0.334 & -0.009 & 0.615 & +0.045 \\
     7 & 0.34  & -0.003 & 0.575 & +0.005 \\
     8 & 0.355 & +0.012 & 0.545 & -0.026 \\
     9 & 0.36  & +0.017 & 0.57  & -0.001 \\
    10 & 0.345 & +0.002 & 0.56  & -0.011 \\
    \midrule
    \multicolumn{1}{@{}l}{Mittelwert $\bar{x}$} & 0.343 & & 0.571 & \\
    \multicolumn{1}{@{}l}{Standardabweichung $s$} & & 0.011 & & 0.021 \\
    \multicolumn{1}{@{}l}{Spanne $R$} & & 0.040 & & 0.070 \\
    \bottomrule
  \end{tabular}
\end{table}

%----------------------------------------------------------------------
% T-02
%----------------------------------------------------------------------
\begin{table}[htbp]
  \centering
  \caption[Messwerte T-02: Vibrationseinheit]{Messwerte des Versuchs T-02 zur
    Auslenkung der Vibrationseinheit in Abhängigkeit der Beladung, nach steigender
    Beladung geordnet. Setup: Anfahren des Endschalters, Vibration, einige Schritte
    höher, anschließend Messung der Abweichung. $x$ bezeichnet den Ablesewert,
    $\Delta$ die Abweichung vom Mittelwert der jeweiligen Messreihe nach
    \cref{eq:abweichung}. Die Masse ist aus der Stückzahl und der protokollierten
    Einzelmasse von \qty{22.9}{\gram} je Hairpin berechnet. Fortsetzung in
    \cref{tab:mess_T02b}.}
  \label{tab:mess_T02a}
  \small
  \setlength{\tabcolsep}{4pt}
  \begin{tabular}{@{}c *{8}{S[table-format=+1.3]}@{}}
    \toprule
     & \multicolumn{2}{c}{\textbf{0}} & \multicolumn{2}{c}{\textbf{10}} & \multicolumn{2}{c}{\textbf{25}} & \multicolumn{2}{c}{\textbf{40}} \\
    \cmidrule(lr){2-3}\cmidrule(lr){4-5}\cmidrule(lr){6-7}\cmidrule(l){8-9}
    \multicolumn{1}{@{}l}{Masse / \unit{\kilo\gram}} & \multicolumn{2}{c}{0.000} & \multicolumn{2}{c}{0.229} & \multicolumn{2}{c}{0.573} & \multicolumn{2}{c}{0.916} \\
    \multicolumn{1}{@{}l}{Normierung / \unit{\milli\meter}} & \multicolumn{2}{c}{0.10} & \multicolumn{2}{c}{0.50} & \multicolumn{2}{c}{0.50} & \multicolumn{2}{c}{0.50} \\
    \midrule
    \textbf{Durchlauf} & {$x$} & {$\Delta$} & {$x$} & {$\Delta$} & {$x$} & {$\Delta$} & {$x$} & {$\Delta$} \\
     & {\unit{\milli\meter}} & {\unit{\milli\meter}} & {\unit{\milli\meter}} & {\unit{\milli\meter}} & {\unit{\milli\meter}} & {\unit{\milli\meter}} & {\unit{\milli\meter}} & {\unit{\milli\meter}} \\
    \midrule
     1 & 0.400 & +0.207 & 0.830 & +0.139 & 0.810 & +0.179 & 0.190 & -0.145 \\
     2 & 0.100 & -0.093 & 0.835 & +0.144 & 0.490 & -0.142 & 0.530 & +0.195 \\
     3 & 0.390 & +0.197 & 0.835 & +0.144 & 0.805 & +0.174 & 0.190 & -0.145 \\
     4 & 0.100 & -0.093 & 0.510 & -0.181 & 0.485 & -0.147 & 0.530 & +0.195 \\
     5 & 0.410 & +0.217 & 0.830 & +0.139 & 0.800 & +0.169 & 0.000 & -0.335 \\
     6 & 0.090 & -0.103 & 0.510 & -0.181 & 0.485 & -0.147 & 0.530 & +0.195 \\
     7 & 0.400 & +0.207 & 0.840 & +0.149 & 0.810 & +0.179 & 0.180 & -0.155 \\
     8 & 0.000 & -0.193 & 0.520 & -0.171 & 0.490 & -0.142 & 0.525 & +0.190 \\
     9 & 0.400 & +0.207 & 0.840 & +0.149 & 0.810 & +0.179 & 0.185 & -0.150 \\
    10 & 0.000 & -0.193 & 0.525 & -0.166 & 0.480 & -0.152 & 0.520 & +0.185 \\
    11 & 0.390 & +0.197 & 0.850 & +0.159 & 0.790 & +0.159 & 0.180 & -0.155 \\
    12 & 0.100 & -0.093 & 0.520 & -0.171 & 0.460 & -0.172 & 0.520 & +0.185 \\
    13 & 0.400 & +0.207 & 0.840 & +0.149 & 0.780 & +0.149 & 0.180 & -0.155 \\
    14 & 0.100 & -0.093 & 0.520 & -0.171 & 0.450 & -0.182 & 0.515 & +0.180 \\
    15 & 0.110 & -0.083 & 0.830 & +0.139 & 0.780 & +0.149 & 0.110 & -0.225 \\
    16 & 0.000 & -0.193 & 0.520 & -0.171 & 0.450 & -0.182 & 0.510 & +0.175 \\
    17 & 0.090 & -0.103 & 0.825 & +0.134 & 0.775 & +0.144 & 0.170 & -0.165 \\
    18 & 0.000 & -0.193 & 0.500 & -0.191 & 0.455 & -0.177 & 0.510 & +0.175 \\
    19 & 0.380 & +0.187 & 0.830 & +0.139 & 0.775 & +0.144 & 0.115 & -0.220 \\
    20 & 0.000 & -0.193 & 0.510 & -0.181 & 0.450 & -0.182 & 0.510 & +0.175 \\
    \midrule
    \multicolumn{1}{@{}l}{Mittelwert $\bar{x}$} & 0.193 &  & 0.691 &  & 0.632 &  & 0.335 &  \\
    \multicolumn{1}{@{}l}{Standardabw. $s$} &  & 0.175 &  & 0.163 &  & 0.167 &  & 0.194 \\
    \multicolumn{1}{@{}l}{Spanne $R$} &  & 0.410 &  & 0.350 &  & 0.360 &  & 0.530 \\
    \bottomrule
  \end{tabular}
\end{table}

\begin{table}[htbp]
  \centering
  \caption[Messwerte T-02 (Fortsetzung)]{Fortsetzung von \cref{tab:mess_T02a}.}
  \label{tab:mess_T02b}
  \small
  \setlength{\tabcolsep}{4pt}
  \begin{tabular}{@{}c *{6}{S[table-format=+1.3]}@{}}
    \toprule
     & \multicolumn{2}{c}{\textbf{60}} & \multicolumn{2}{c}{\textbf{80}} & \multicolumn{2}{c}{\textbf{100}} \\
    \cmidrule(lr){2-3}\cmidrule(lr){4-5}\cmidrule(l){6-7}
    \multicolumn{1}{@{}l}{Masse / \unit{\kilo\gram}} & \multicolumn{2}{c}{1.374} & \multicolumn{2}{c}{1.832} & \multicolumn{2}{c}{2.290} \\
    \multicolumn{1}{@{}l}{Normierung / \unit{\milli\meter}} & \multicolumn{2}{c}{3.50} & \multicolumn{2}{c}{3.50} & \multicolumn{2}{c}{3.50} \\
    \midrule
    \textbf{Durchlauf} & {$x$} & {$\Delta$} & {$x$} & {$\Delta$} & {$x$} & {$\Delta$} \\
     & {\unit{\milli\meter}} & {\unit{\milli\meter}} & {\unit{\milli\meter}} & {\unit{\milli\meter}} & {\unit{\milli\meter}} & {\unit{\milli\meter}} \\
    \midrule
     1 & 3.120 & +0.128 & 3.205 & -0.145 & 1.100 & -0.148 \\
     2 & 3.110 & +0.118 & 3.500 & +0.151 & 0.000 & -1.248 \\
     3 & 3.120 & +0.128 & 3.210 & -0.140 & 0.000 & -1.248 \\
     4 & 2.770 & -0.222 & 3.490 & +0.141 & 0.000 & -1.248 \\
     5 & 3.420 & +0.428 & 3.200 & -0.150 & 0.000 & -1.248 \\
     6 & 2.790 & -0.202 & 3.500 & +0.151 & 0.000 & -1.248 \\
     7 & 3.100 & +0.108 & 3.205 & -0.145 & 2.610 & +1.362 \\
     8 & 2.780 & -0.212 & 3.500 & +0.151 & 0.000 & -1.248 \\
     9 & 3.430 & +0.438 & 3.205 & -0.145 & 0.000 & -1.248 \\
    10 & 2.785 & -0.207 & 3.500 & +0.151 & 2.270 & +1.022 \\
    11 & 3.110 & +0.118 & 3.205 & -0.145 & 0.000 & -1.248 \\
    12 & 2.790 & -0.202 & 3.500 & +0.151 & 3.300 & +2.052 \\
    13 & 3.100 & +0.108 & 3.205 & -0.145 & 0.000 & -1.248 \\
    14 & 2.780 & -0.212 & 3.495 & +0.146 & 0.000 & -1.248 \\
    15 & 3.100 & +0.108 & 3.200 & -0.150 & 0.000 & -1.248 \\
    16 & 2.780 & -0.212 & 3.495 & +0.146 & 7.700 & +6.452 \\
    17 & 3.090 & +0.098 & 3.200 & -0.150 & 0.000 & -1.248 \\
    18 & 2.780 & -0.212 & 3.495 & +0.146 & 4.810 & +3.562 \\
    19 & 3.100 & +0.108 & 3.200 & -0.150 & 0.000 & -1.248 \\
    20 & 2.780 & -0.212 & 3.480 & +0.131 & 3.170 & +1.922 \\
    \midrule
    \multicolumn{1}{@{}l}{Mittelwert $\bar{x}$} & 2.992 &  & 3.350 &  & 1.248 &  \\
    \multicolumn{1}{@{}l}{Standardabw. $s$} &  & 0.216 &  & 0.150 &  & 2.120 \\
    \multicolumn{1}{@{}l}{Spanne $R$} &  & 0.660 &  & 0.300 &  & 7.700 \\
    \bottomrule
  \end{tabular}
\end{table}

Die in \cref{tab:mess_T02a,tab:mess_T02b} angegebenen Kennwerte $s$ und $R$ beziehen
sich auf die jeweils vollständige Messreihe. Sie sind nur eingeschränkt als
Streuungsmaß zu lesen, da die Ablesewerte innerhalb einer Messreihe zwischen zwei
wiederkehrenden Niveaus wechseln und der Mittelwert damit zwischen beiden Niveaus zu
liegen kommt. Dies zeigt sich auch im Vorzeichenwechsel der Spalte $\Delta$ von
Durchlauf zu Durchlauf. Eine getrennte Auswertung der beiden Niveaus findet sich in
\cref{tab:mess_T02_zustaende}.

\begin{table}[htbp]
  \centering
  \caption[T-02: Auswertung nach Niveaus]{Getrennte Auswertung der beiden
    wiederkehrenden Ablesenivaus aus \cref{tab:mess_T02a,tab:mess_T02b}. Die Zuordnung erfolgt
    über den Median der jeweiligen Messreihe. Die Amplitude bezeichnet den Abstand
    der beiden Niveaumittelwerte und ist als einzige Kenngröße dieses Versuchs
    unabhängig vom Normierungswert.}
  \label{tab:mess_T02_zustaende}
  \small
  \begin{tabular}{@{}r S[table-format=1.3] S[table-format=1.3] S[table-format=1.3] S[table-format=1.3] S[table-format=1.3]@{}}
    \toprule
    \textbf{Beladung} & \multicolumn{2}{c}{\textbf{unteres Niveau}}
      & \multicolumn{2}{c}{\textbf{oberes Niveau}} & {\textbf{Amplitude}} \\
    \cmidrule(lr){2-3}\cmidrule(lr){4-5}\cmidrule(l){6-6}
    {Hairpins} & {$\bar{x}$} & {$s$} & {$\bar{x}$} & {$s$} & {$\Delta\bar{x}$} \\
    & {\unit{\milli\meter}} & {\unit{\milli\meter}} & {\unit{\milli\meter}}
    & {\unit{\milli\meter}} & {\unit{\milli\meter}} \\
    \midrule
      0 & 0.053 & 0.051 & 0.364 & 0.096 & 0.312 \\
     10 & 0.546 & 0.098 & 0.836 & 0.007 & 0.290 \\
     25 & 0.470 & 0.018 & 0.794 & 0.015 & 0.324 \\
     40 & 0.150 & 0.060 & 0.520 & 0.008 & 0.370 \\
     60 & 2.812 & 0.098 & 3.171 & 0.134 & 0.359 \\
     80 & 3.204 & 0.003 & 3.495 & 0.006 & 0.292 \\
    100 & 0.000 & 0.000 & 3.566 & 2.142 & 3.566 \\
    \bottomrule
  \end{tabular}
\end{table}

Für die Beladungen von 10 bis 80 Hairpins liegt die Amplitude zwischen
\qty{0.290}{\milli\meter} und \qty{0.370}{\milli\meter} und zeigt damit keine
ausgeprägte Abhängigkeit von der Beladung. Die Messreihe mit 100 Hairpins fällt aus
diesem Bild heraus: 13 der 20 Ablesungen betragen exakt null, die übrigen streuen mit
$s = \qty{2.142}{\milli\meter}$ ohne erkennbare Systematik. Bei einer Beladung von
\qty{2.29}{\kilo\gram} führt die Vibrationseinheit die Bewegung folglich in der
Mehrzahl der Durchläufe nicht mehr aus.

%----------------------------------------------------------------------
% T-03
%----------------------------------------------------------------------
\begin{table}[htbp]
  \centering
  \caption[Messwerte T-03: Transportsystem]{Messwerte des Versuchs T-03 zur
    Positionsstreuung des Transportsystems im open loop. Setup: Endstop anfahren,
    3000 Schritte zurück und halten, 200 Schritte hin und her, 300 Schritte extra
    für die Messung. Normierung auf 350.}
  \label{tab:mess_T03}
  \small
  \begin{tabular}{@{}cS[table-format=3.1]cS[table-format=3.1]@{}}
    \toprule
    \textbf{RUN} & {\textbf{Abweichung}} & \textbf{RUN} & {\textbf{Abweichung}} \\
    \midrule
     1 & 359   & 11 & 358.5 \\
     2 & 358   & 12 & 358.5 \\
     3 & 358.5 & 13 & 362   \\
     4 & 358.5 & 14 & 359   \\
     5 & 359   & 15 & 359   \\
     6 & 358   & 16 & 364   \\
     7 & 358.5 & 17 & 358.5 \\
     8 & 360.5 & 18 & 359   \\
     9 & 360   & 19 & 359   \\
    10 & 358   & 20 & 360   \\
    \bottomrule
  \end{tabular}
\end{table}

%----------------------------------------------------------------------
% T-04
%----------------------------------------------------------------------
\begin{table}[htbp]
  \centering
  \caption[Messwerte T-04: Fixierungssystem]{Messwerte des Versuchs T-04 zur
    Lagetreue des Hairpins in der Fixierung. Setup: Teilung in zwei Richtungen,
    reines Aufstellen und seitliche Verschiebung. Normierung auf 50.}
  \label{tab:mess_T04}
  \small
  \begin{tabularx}{\textwidth}{@{}cS[table-format=2.1]X@{}}
    \toprule
    \textbf{RUN} & {\textbf{Abweichung}} & \textbf{Bemerkung} \\
    \midrule
     1 & 21   & Normierung auf 50 \\
     2 & 71.5 & Hairpin wird im Trichter gelassen \\
     3 & 61.5 & Servos schalten voll durch \\
     4 & 64.5 & keine Vibration der Linearachse \\
     5 & 60   & \\
     6 & 49   & \\
     7 & 39.5 & \\
     8 & 71   & \\
     9 & 40   & \\
    10 & 63   & \\
    11 & 64   & \\
    12 & 31.5 & \\
    13 & 47   & \\
    14 & 45   & \\
    15 & 62   & \\
    16 & 20   & \\
    17 & 19   & \\
    18 & 61   & \\
    19 & 56   & \\
    20 & 49.5 & \\
    \bottomrule
  \end{tabularx}
\end{table}

%----------------------------------------------------------------------
% T-05
%----------------------------------------------------------------------
\begin{table}[htbp]
  \centering
  \caption[Geometrien der Hairpins in Versuch T-05]{Geometrien der in Versuch T-05
    verwendeten Hairpin-Varianten. Alle Durchläufe wurden mit dem Endeffektor
    TestV1 durchgeführt.}
  \label{tab:mess_T05_geom}
  \small
  \begin{tabularx}{\textwidth}{@{}l S[table-format=2.2] S[table-format=1.2] S[table-format=3.0] X@{}}
    \toprule
    \textbf{Variante} & {\textbf{Breite}} & {\textbf{Kupferbreite}}
      & {\textbf{Länge}} & \textbf{Bemerkung} \\
    & {\unit{\milli\meter}} & {\unit{\milli\meter}} & {\unit{\milli\meter}} & \\
    \midrule
    A & 72.74 & 2.15 & 155 & nicht entgratete Füße, oberste Lage Wafios, kupferfarben \\
    B & 65    & 2.15 & 153 & verschieden lange Füße, zweite Lage von oben, Wafios \\
    C & 64.5  & 2.3  & 145 & fünfte Lage von oben, großer Ständer \\
    D & 69.55 & 2.30 & 184 & zweite Lage von oben, großer Ständer \\
    \bottomrule
  \end{tabularx}
\end{table}

\begin{table}[htbp]
  \centering
  \caption[Fehlercodes des Versuchs T-05]{Fehlercodes zur Bewertung der
    Vereinzelung in Versuch T-05.}
  \label{tab:mess_T05_codes}
  \small
  \begin{tabularx}{\textwidth}{@{}cX@{}}
    \toprule
    \textbf{Code} & \textbf{Bedeutung} \\
    \midrule
    0 & kein Fehler, alles innerhalb der Toleranz \\
    1 & schräg im Trichter (seitwärts) \\
    2 & doppelte Ausgabe \\
    3 & Herausspringen aus der Rutsche \\
    4 & schräg aus dem Trichter gefallen \\
    5 & ein Bein im Trichter, eins außerhalb \\
    \bottomrule
  \end{tabularx}
\end{table}

\begin{table}[htbp]
  \centering
  \caption[Messwerte T-05: Vereinzelung]{Messwerte des Versuchs T-05 zur
    Vereinzelungstreue. Setup: 25 Hairpins geladen, Homing-Vibration, Vereinzeln,
    \qty{10}{\second} warten, anschließend Entnahme des Hairpins aus dem Hopper.
    Die Bewertung erfolgt nach den Fehlercodes aus \cref{tab:mess_T05_codes},
    die Varianten A bis D sind in \cref{tab:mess_T05_geom} beschrieben.}
  \label{tab:mess_T05}
  \small
  \begin{tabular}{@{}ccccc@{}}
    \toprule
    & \multicolumn{4}{c}{\textbf{Fehlercode je Hairpin-Variante}} \\
    \cmidrule(l){2-5}
    \textbf{RUN} & \textbf{A} & \textbf{B} & \textbf{C} & \textbf{D} \\
    \midrule
     1 & 0 & 2 & 0 & 0 \\
     2 & 3 & 2 & 0 & 0 \\
     3 & 3 & 0 & 0 & 0 \\
     4 & 0 & 2 & 0 & 0 \\
     5 & 0 & 2 & 0 & 0 \\
     6 & 0 & 0 & 0 & 0 \\
     7 & 0 & 2 & 0 & 0 \\
     8 & 1 & 0 & 0 & 0 \\
     9 & 0 & 2 & 0 & 0 \\
    10 & 0 & 0 & 0 & 0 \\
    11 & 0 & 2 & 0 & 0 \\
    12 & 0 & 2 & 0 & 0 \\
    13 & 0 & 2 & 0 & 0 \\
    14 & 1 & 2 & 0 & 0 \\
    15 & 3 & 2 & 5 & 0 \\
    16 & 4 & 2 & 0 & 0 \\
    17 & 4 & 2 & 4 & 2 \\
    18 & 3 & 2 & 0 & 2 \\
    19 & 0 & 2 & 0 & 2 \\
    20 & 5 & 2 & 0 & 0 \\
    21 & 5 & 2 & 0 & 0 \\
    22 & 0 & 0 & 0 & 0 \\
    23 & 0 & 0 & 0 & 0 \\
    24 & 0 & 2 & 0 & 2 \\
    25 & 0 & 0 & 0 & 0 \\
    \bottomrule
  \end{tabular}
\end{table}

%----------------------------------------------------------------------
% T-06
%----------------------------------------------------------------------
\begin{table}[htbp]
  \centering
  \caption[Messwerte T-06: Servomotoren des Fixierungssystems]{Messwerte des
    Versuchs T-06 zu den Servomotoren des Fixierungssystems. Setup: 20 Durchläufe
    je Seite, Normierung auf 350. Bei Servo~2 konnte die Messuhr nur schlecht
    platziert werden, weshalb dort mit einer erhöhten Messunsicherheit zu rechnen
    ist. Für Servo~3 war die Platzierung besser.}
  \label{tab:mess_T06}
  \small
  \begin{tabular}{@{}cS[table-format=3.1]S[table-format=3.1]@{}}
    \toprule
    \textbf{RUN} & {\textbf{Servo 2}} & {\textbf{Servo 3}} \\
    \midrule
     1 & 176   & 343   \\
     2 & 172   & 342   \\
     3 & 176   & 343   \\
     4 & 176.5 & 340   \\
     5 & 176   & 339.5 \\
     6 & 175   & 337   \\
     7 & 176   & 336.5 \\
     8 & 176   & 333   \\
     9 & 176   & 331   \\
    10 & 176   & 329   \\
    11 & 176   & 329   \\
    12 & 175.5 & 329   \\
    13 & 176.5 & 328   \\
    14 & 176   & 327   \\
    15 & 175   & 324   \\
    16 & 175.5 & 325   \\
    17 & 175   & 324   \\
    18 & 175.5 & 321   \\
    19 & 175   & 321   \\
    20 & 175.5 & 321   \\
    \bottomrule
  \end{tabular}
\end{table}

%----------------------------------------------------------------------
% T-07
%----------------------------------------------------------------------
\begin{table}[htbp]
  \centering
  \caption[Messwerte T-07: Positionierung des Transportsystems]{Messwerte des
    Versuchs T-07 zur Positionierung im Gesamtverbund des Transportsystems.
    Setup: 25 Hairpins unterschiedlicher Länge, Homen der Linearachse, Anfahren
    der Position, Ablesen der Abweichung mit dem Messschieber. Durchgang~1 wurde
    mit einem \qty{64}{\milli\meter}-Pin aus der fünften Reihe des Ständers
    durchgeführt (Normierung längs \qty{5}{\milli\meter}, schräg
    \qty{4}{\milli\meter}), Durchgang~2 mit einem \qty{69}{\milli\meter}-Pin aus
    der zweiten Reihe (Normierung längs \qty{7.5}{\milli\meter}, schräg
    \qty{4}{\milli\meter}).}
  \label{tab:mess_T07}
  \small
  \begin{tabular}{@{}c S[table-format=1.2] S[table-format=1.2] S[table-format=1.2] S[table-format=1.2]@{}}
    \toprule
    & \multicolumn{2}{c}{\textbf{Durchgang 1 (\qty{64}{\milli\meter})}}
    & \multicolumn{2}{c}{\textbf{Durchgang 2 (\qty{69}{\milli\meter})}} \\
    \cmidrule(lr){2-3}\cmidrule(l){4-5}
    \textbf{RUN} & {\textbf{schräg}} & {\textbf{längs}}
                 & {\textbf{schräg}} & {\textbf{längs}} \\
    & {\unit{\milli\meter}} & {\unit{\milli\meter}}
    & {\unit{\milli\meter}} & {\unit{\milli\meter}} \\
    \midrule
     1 & 7.33 & 2.22 & 7.4  & 3.78 \\
     2 & 7.93 & 2.11 & 7.98 & 3.86 \\
     3 & 7.52 & 1.85 & 7.63 & 4.65 \\
     4 & 7.25 & 1.92 & 8.01 & 4.94 \\
     5 & 7.19 & 1.53 & 7.88 & 4.49 \\
     6 & 8.09 & 1.85 & 7.81 & 4.79 \\
     7 & 7.59 & 1.81 & 8.1  & 4.66 \\
     8 & 7.54 & 1.45 & 7.63 & 5.4  \\
     9 & 7.97 & 1.42 & 8.13 & 4.54 \\
    10 & 7.7  & 2.05 & 8.16 & 4.4  \\
    11 & 7.85 & 1.63 & 7.99 & 4.99 \\
    12 & 8.1  & 1.6  & 8.15 & 5.07 \\
    13 & 7.53 & 1.36 & 8.18 & 4.71 \\
    14 & 7.9  & 1.69 & 8.1  & 5.14 \\
    15 & 7.74 & 1.83 & 7.79 & 4.91 \\
    16 & 7.89 & 1.66 & 8.24 & 5.3  \\
    17 & 7.99 & 1.91 & 8.19 & 5.2  \\
    18 & 7.35 & 1.61 & 8.24 & 5.21 \\
    19 & 8.25 & 2.26 & 7.89 & 5.5  \\
    20 & 7.61 & 1.24 & 7.82 & 4.82 \\
    21 & 7.99 & 1.9  & 7.99 & 5.08 \\
    22 & 7.8  & 1.72 & 8.14 & 5.58 \\
    23 & 7.48 & 1.66 & 7.94 & 5.31 \\
    24 & 7.57 & 1.82 & 7.69 & 5.09 \\
    25 & 8.24 & 2.1  & 8.1  & 5.5  \\
    \bottomrule
  \end{tabular}
\end{table}

Zusätzlich zu den in \cref{tab:mess_T07} aufgeführten Abweichungen wurde in Versuch
T-07 die Höhenlage geprüft. Dabei wurde der \qty{69}{\milli\meter}-Pin aus der
zweiten Reihe des Ständers bei einer Normierung auf \qty{0}{\milli\meter}
verwendet. Alle 25 Durchläufe wurden ohne feststellbare Höhenabweichung als
einwandfrei bewertet.

%----------------------------------------------------------------------
% T-08
%----------------------------------------------------------------------
\begin{table}[htbp]
  \centering
  \caption[Messwerte T-08: Gesamtsystem]{Ergebnisse des Versuchs T-08 zum
    Gesamtdurchlauf bei einer Taktzeit von \qty{43}{\second}. Je Durchgang wurden
    zehn Durchläufe protokolliert und nach fehlerfreiem Ablauf beziehungsweise
    aufgetretener Störung bewertet.}
  \label{tab:mess_T08}
  \small
  \begin{tabularx}{\textwidth}{@{}cXXX@{}}
    \toprule
    \textbf{RUN} & \textbf{Durchgang 1} & \textbf{Durchgang 2}
      & \textbf{Durchgang 3} \\
    \midrule
     1 & fehlerfrei & fehlerfrei & fehlerfrei \\
     2 & fehlerfrei & fehlerfrei & fehlerfrei \\
     3 & fehlerfrei & fehlerfrei & Vereinzelung verhakt \\
     4 & Timing     & fehlerfrei & fehlerfrei \\
     5 & fehlerfrei & fehlerfrei & fehlerfrei \\
     6 & fehlerfrei & fehlerfrei & fehlerfrei \\
     7 & fehlerfrei & Vereinzelung verhakt & fehlerfrei \\
     8 & Timing     & fehlerfrei & fehlerfrei \\
     9 & fehlerfrei & fehlerfrei & fehlerfrei \\
    10 & fehlerfrei & fehlerfrei & Timing \\
    \bottomrule
  \end{tabularx}
\end{table}

\clearpage
